% Owner: Tân. Numbers come from Chapter 5 (output/inference/).
\section{Main Findings}
\begin{itemize}
\item \textbf{Width is the main size signal.} Ads are wider than non-ads by
  about 201~px on average (95\% CI 181 to 222~px). With height, Local and
  missingness held fixed, each doubling of width multiplies the odds of an
  ad by about 5.2. There is no evidence of a height association once width
  is known.
\item \textbf{Missing size information is a weak signal.} Records with a
  missing feature are less often ads (11.0\% against 17.3\%), but the
  association is weak (Cram\'er's $V=0.071$) and is not significant after
  adjustment for size and Local.
\item \textbf{The words around an image carry the most information.} The
  strongest terms read like advertising: \texttt{url*ads},
  \texttt{ancurl*click}, \texttt{ancurl*com} and \texttt{alt*click}. Adding
  six terms to the main-effects model raises test average precision from
  0.594 to 0.795 (gain 0.200, 95\% CI 0.092 to 0.318) and catches 23 more
  of the 77 test ads (65 against 42), at the cost of 20 more false alarms.
\item \textbf{A curved size effect improves ranking, not decisions.} The
  format + interaction model raises test average precision by 0.090 (95\%
  CI 0.039 to 0.145), but at the chosen threshold it catches only 3 more
  ads, a difference consistent with chance.
\end{itemize}

\section{Strengths of the Methods}
\begin{itemize}
\item \textbf{An honest test set.} Exact duplicates were removed before the
  analysis, and records with identical features were kept on one side of
  the split. The test set therefore measures performance on records the
  models have not seen.
\item \textbf{No information from the test set.} Missing-value filling,
  term selection and decision thresholds were all set on training data,
  inside each cross-validation fold where needed, and then frozen.
\item \textbf{Interpretable models.} Logistic regression gives odds ratios
  with confidence intervals, so each conclusion states both a direction and
  a size.
\item \textbf{Uncertainty is shown, not hidden.} Confidence intervals,
  McNemar tests and bootstrap intervals separate the gains that are clearly
  real (the terms) from those that may be chance (the extra ads caught by
  the curved size effect).
\item \textbf{Missingness kept as information.} The \texttt{missing\_size}
  flag allowed us to test whether missing values relate to ad status,
  instead of hiding them by filling them in.
\end{itemize}

\section{Limitations}
\begin{itemize}
\item \textbf{One small test set.} All test results rest on one split with
  77 ads, so intervals are wide. A different split could change the
  smaller differences.
\item \textbf{Terms chosen by the data.} The six terms were selected because
  they differ most between classes, and the same indicators had been seen in
  the exploratory analysis. Their tests describe strong associations in this
  dataset; they do not confirm a hypothesis fixed in advance.
\item \textbf{Overlapping terms.} The three \texttt{alt} terms nearly
  always occur together, so their separate odds ratios are imprecise, and
  some term combinations occur almost only among ads, which makes the odds
  ratios very large and unstable.
\item \textbf{The precision target was not kept on new records.} The term
  model reached a training precision of 0.70, but its test precision
  interval, [0.51, 0.69], lies just below that target.
\item \textbf{Independence and representativeness.} The tests assume
  independent records, but several images can come from the same web page
  or site, and only exact duplicates were removed. The data come from a
  1998 collection, and advertising formats have changed since then.
\item \textbf{Observed widths only.} The width test uses records with an
  observed width, and missingness differs between the classes, so the
  tested records do not fully represent the images with unknown size.
\end{itemize}

\section{Suggested Extensions}
\begin{itemize}
\item \textbf{Confirm a few terms chosen in advance.} Fix a small set of
  terms by reasoning before looking at the data (for example ``ads'' and
  ``click'') and test them on new data, so that the result is confirmatory.
\item \textbf{Use all terms with penalized regression.} Lasso logistic
  regression can consider all 1,554 terms at once and handles overlapping
  terms by shrinking their coefficients.
\item \textbf{Repeat the split.} Repeated grouped cross-validation would
  show how much the test results vary between splits, instead of relying on
  one split.
\item \textbf{Set the threshold by costs.} In practice, the threshold should
  reflect the relative cost of a missed ad and of a false alarm, and it
  should be checked again on recent data from other websites.
\end{itemize}
